% ---------------------------------------------------------------- 第一部分导言
\section*{第一部分导语：周期视角下的中国农业}
\addcontentsline{toc}{section}{第一部分导语}

分析农业周期最常见的方法，是把农产品价格当作一条宏观时间序列，
用 CPI 食品分项或农产品批发价格指数去解释它的波动。
本部分不采用这一路径，理由是：\hl{农业的周期不是一条曲线，而是许多条互不同步的曲线}。
玉米与生猪的价格可以在同一年分别上涨与下跌，
苹果与红枣的价格可以在同一年相差数十个百分点（第 \ref{sec:soft} 章），
牛的价格周期长达四年而蛋鸡只有半年（第 \ref{sec:other-livestock} 章）。

因此第一部分的结构安排是：
\begin{enumerate}[leftmargin=1.5em, itemsep=2pt]
  \item \textbf{先建立分析框架}（第 \ref{sec:framework} 章）：
        用申万行业分类（2021 版）确定讨论边界，
        把需求拆成“消费、工业原料、饲料”三条链，
        把供给拆成“当期产量”与“滞后产能”，
        并说明本书的数据口径与统一计量方法。
  \item \textbf{再逐品种建立周期模型}：生猪（第 \ref{sec:hog} 章，基准周期）、
        粮食与油料（第 \ref{sec:grain} 章）、
        软商品与经济作物（第 \ref{sec:soft} 章）、
        其他养殖与水产（第 \ref{sec:other-livestock} 章）。
        每一章都回答三个问题：周期从哪里来？现在处在哪一段？什么会终止它？
  \item \textbf{最后做跨品种检验}（第 \ref{sec:corr} 章）：
        以猪周期为基准，在\hl{同一时间窗口、同一频率}上计算
        相关矩阵、领先滞后关系与传导回归，
        并检验这种相关性是否延伸到股票市场（申万各子行业指数）。
\end{enumerate}

贯穿全书的三个方法论约定，需要在开篇明确：

\noindent\textbf{① 相关性只在同一窗口、同一频率上比较。}
不同品种的期货上市时间差异极大（豆粕自 \num{2005} 年、生猪自 \num{2021} 年、
20 号胶自 \num{2019} 年），若各自使用最长样本，
得到的相关系数反映的是不同的宏观环境。本书统一使用
\num{2021} 年 2 月---\num{2026} 年 8 月的月度窗口，
并以更长样本（\num{2015} 年以来）做稳健性检验。

\noindent\textbf{② 周期性用“锯齿（ZigZag）”划分，而不是用主观描述。}
定义价格从局部极值反向变动超过阈值（生猪期货 \num{18}\%、
长样本现货 \num{20}\%）为周期转换，从而把价格序列机械地切分为
主升段与主跌段，再统计各段的持续时间与涨跌幅。

\noindent\textbf{③ 所有出现在正文的数字都由脚本从原始数据生成。}
报告仓库包含完整的抓取、清洗、计算与制表脚本，
正文表格通过 \texttt{\textbackslash input} 引用自动生成的 \LaTeX{} 文件，
\hl{不存在手工誊抄的数据}。
\clearpage
